\documentclass[../main.tex]{subfiles}

\begin{document}

\chapter{Brief Dataset Introduction}
\label{chap:dataset}


% ============================================================
\section{Introduction to the CIC-IDS2018 Dataset}
\label{sec:dataset-intro}
% ============================================================

The dataset used in the experimental phase of this work is the \textbf{CIC-IDS2018}, a publicly available benchmark released by the University of New Brunswick.  It was produced as part of a systematic effort to build realistic, labelled network-traffic traces that can drive the development and evaluation of next-generation intrusion-detection systems (IDS).

The data was collected over a five day period in February–March 2018 inside a controlled cyber-range infrastructure.  A set of attack machines (an external network of 50 machines) launched a curated battery of attacks against a victim network of 420 machines and 30 servers, while the defender-side profile generator replayed realistic benign traffic drawn from a profiled usage model that incorporates HTTP, HTTPS, FTP, SSH, and email activity. All bi-directional network flows were captured with the \textit{CIC\-Flow\-Meter} tool\supercite{cicflowmeter}, which extracts statistical features: primarily packet-level and flow-level statistics from raw \texttt{.pcap} dumps. Each flow is labelled with either \texttt{Benign} or one of fourteen mutually exclusive attack categories.

The attack repertoire covers a broad spectrum of modern threat types:
\begin{itemize}
    \item \textbf{Brute-Force} attacks (FTP-BruteForce, SSH-BruteForce);
    \item \textbf{Denial-of-Service (DoS)} attacks (DoS-GoldenEye, DoS-Hulk,\newline DoS-SlowHTTPTest, DoS-Slowloris);
    \item \textbf{Web attacks} (SQL Injection, XSS, Brute-Force on web login);
    \item \textbf{Botnet} traffic (ARES botnet);
    \item \textbf{Infiltration} scenarios;
    \item \textbf{Distributed Denial-of-Service (DDoS)} attacks with LOIC UDP and HOIC.
\end{itemize}

The dataset was originally distributed as a collection of day-labelled CSV files, one per day of capture.  A pre-processing step prior to the analysis merged these individual files into a unified dataset and resolved a minor schema inconsistency: four features present only in one daily capture were not shared across all files; rather than discarding those flows, those anomalous features were saved separately so that no information was irrecoverably lost.
The unified dataset was subsequently persisted in the Apache Parquet columnar format to enable efficient distributed reads via Apache Spark in the experimental pipeline.

The dataset is available for download from the Kaggle platform \supercite{ids2018}, while additional documentation and the original raw captures are hosted at CICIDS offical website\supercite{cicids2018Official}.

% ============================================================
\newpage
\section{Dimensionality and Schema}
\label{sec:dataset-dimensionality}
% ============================================================

After unification, the dataset contains \textbf{16\,173\,208 network-flow instances} described by \textbf{79 features} plus two label columns: \texttt{Label} and \texttt{Label\_generic}).
The 79 features are purely statistical descriptors extracted by CICFlowMeter;
they include, among others:

\begin{itemize}
  \item \emph{Flow duration and inter-arrival times} (\texttt{Flow Duration},
        \texttt{Flow IAT Mean}, \texttt{Flow IAT Std}, etc.);
  \item \emph{Packet-length statistics} (\texttt{Fwd Pkt Len Mean},
        \texttt{Bwd Pkt Len Max}, \texttt{Pkt Len Var}, etc.);
  \item \emph{Byte and packet counts}, split by direction (\texttt{Tot Fwd
        Pkts}, \texttt{Tot Bwd Pkts}, \texttt{TotLen Fwd Pkts}, etc.);
  \item \emph{TCP flag counts} (\texttt{SYN Flag Cnt}, \texttt{FIN Flag
        Cnt}, \texttt{PSH Flag Cnt}, etc.);
  \item \emph{Rate-based features} (\texttt{Flow Byts/s}, \texttt{Flow
        Pkts/s}, \texttt{Bwd Pkts/s}, etc.);
  \item \emph{Sub-flow and bulk statistics} (\texttt{Subflow Fwd Pkts},
        \texttt{Fwd Blk Rate Avg}, etc.);
  \item \emph{Window size and header-length fields}
        (\texttt{Init\_Win\_bytes\_forward}, \texttt{Fwd Header Len}, etc.).
\end{itemize}

Crucially, the dataset contains \emph{no application-layer payload information}: all features are derived solely from transport- and network-layer flow metadata.  This design choice, motivated by privacy regulations that preclude deep-packet inspection on real traffic, constitutes an inherent upper bound on the discriminative power available to any classifier trained on these data, a limitation that will be revisited when analyzing feature–target correlations (Chapter \ref{sec:feat-target-corr}).

% ============================================================
\newpage
\section{Label Distribution and Class Imbalance}
\label{sec:label-distribution}
% ============================================================

A preliminary examination of the label columns reveals a severe class imbalance, a characteristic endemic to real-world network-traffic datasets. The label schema is twofold:

\begin{enumerate}
  \item \texttt{Label}: the original fine-grained label carrying the specific attack category or \texttt{Benign}.
  \item \texttt{Label\_generic}: a binary encoding that we manually inserted in the dataset. This was derived from having\texttt{Label}: $0$ for benign flows and $1$ for any attack flow, regardless of type.  This binarised label was introduced to simplify the experimental evaluation of deep learning models in scenarios where the primary objective is anomaly detection rather than attack classification.
\end{enumerate}

Figure~\ref{fig:label-distribution} shows the joint distribution of both label
hierarchies.  The left pie chart illustrates the overall fraction of benign
versus malign traffic, while the right chart details the relative proportions
of the individual attack categories within the malign subset.

\begin{figure}[htbp]
  \centering
  \includegraphics[width=\textwidth]{images/label_distribution_piecharts.png}
  \caption{Class distribution of the CIC-IDS2018 dataset.  \emph{Left:}
    proportion of benign versus malign flows (binary label).  \emph{Right:}
    breakdown of specific attack categories within the malign portion.}
  \label{fig:label-distribution}
\end{figure}

We can denote that the imbalance is substantial: benign traffic constitutes the overwhelming
majority of the dataset, while the aggregate of all attack types represents a
considerably smaller fraction.  Within the attack classes, further heterogeneity
is observed: volumetric attacks such as DoS/DDoS generate large numbers of
flows and are therefore over-represented relative to low-rate, stealthy attacks
such as infiltration or web-based brute force.

Class imbalance has well-documented implications for machine learning model
training: standard cross-entropy minimization tends to bias models towards the majority class, artificially inflating accuracy while degrading recall on minority attack categories.  Awareness of this imbalance was therefore a prerequisite for designing the experimental protocol, and it also motivates the particular interest in \emph{noise injection} experiments: the introduction of synthetic corruptions by \texttt{PuckTrick} may interact non-uniformly with class-stratified data, potentially exacerbating or alleviating the imbalance in ways that influence classifier behaviour in non-trivial ways.

% ============================================================
\newpage
\section{Temporal Distribution of Traffic}
\label{sec:temporal-distribution}
% ============================================================

Each flow record includes a \texttt{Timestamp} column that encodes the start time of the flow as a character string. Since one of our main interests is to use this dataset as time series a thorough inspection of this column was overdue. This inspection revealed that multiple date-format conventions coexist across daily files (\texttt{yyyy-MM-dd HH:mm:ss}, \texttt{dd/MM/yyyy HH:mm:ss}, etc...).  Additionally, a non-negligible fraction of rows contained epoch-origin timestamps (date 1 January 1970), a classical artefact of integer overflow or uninitialized timestamp fields.  These corrupted entries were filtered out before temporal analysis was performed.

Figure~\ref{fig:temporal-distribution} presents the temporal structure of the dataset along two complementary dimensions.

\begin{figure}[H]
  \centering
  \includegraphics[width=\textwidth]{images/temporal_distribution_label_generic.png}
  \caption{Temporal distribution of network flows.  \emph{Top:}
    stacked area chart showing the hourly volume of benign (blue) and
    malign (red) flows over the entire five-day collection window.
    \emph{Bottom:} aggregated distribution of flow counts by hour of the day
    (0-23), with separate bars for benign and malign traffic.}
  \label{fig:temporal-distribution}
\end{figure}

The stacked area chart (top panel) reveals that attacks are not uniformly spread across the collection period but are concentrated in specific temporal windows. This clustering reflects the scripted nature of the attack sessions: attack tools were activated according to a defined schedule, generating bursts of malign traffic interspersed with periods of exclusively benign activity. This proves to us the validity o using this dataset as a time series: having most of out attack types being of DOS and DDSOS variants we have a lost of subsequent malign packages all compacted in a short period of time. 

The hourly aggregation chart (bottom panel) provides a time-of-day analysis.
Benign traffic follows a circadian pattern consistent with daytime office usage, peaking during working hours, while malign traffic exhibits a more uniform or counter-cyclical distribution depending on the attack type.  

Both of these temporal patterns imply that the \texttt{Timestamp} feature carries genuine discriminative information: a classifier with access to flow timing could leverage temporal context to improve attack detection.

From the standpoint of PuckTrick's noise injection experiments, this observation is significant: Network attacks are not isolated events: they unfold as campaigns with characteristic temporal signatures (mostly as bursty for DDoS and Dos flood).  Models that exploit this temporal context (such as the CNN-LSTM employed in this work) gain a significant advantage, but their temporal reasoning is also vulnerable to corruption of the timestamp feature or of the features that encode temporal dynamics.

% ============================================================
\newpage
\section{Missing and \texttt{inf} Data Analysis}
\label{sec:missing-data}
% ============================================================

Prior to any modelling activity, the completeness of the feature matrix was assessed by counting null, \texttt{NaN}, and \texttt{inf} values across all 79 columns. The dataset scan revealed that one feature contained a small proportion of missing values while another 2 features contained another small fraction of \texttt{NaN} values, for a total of roughly 0.37\% of all rows.

Given that 0.37\% represents a negligible fraction of a dataset containing over 16 million instances, the standard remediation strategy of \emph{listwise deletion} (discarding all rows that contain at least one null value) was adopted.  This approach is justified because:

\begin{enumerate}
  \item The lost rows amount to approximately 60\,000 instances, less than 0.4\% of the total, which is statistically negligible.
  \item The missingness pattern showed no systematic relationship with the target variable; the null-bearing flows were distributed across both benign and attack classes, indicating a process closer to \emph{Missing Completely At Random} (MCAR) rather than a class-specific artefact.
  \item Imputation strategies (mean, median, or model-based) would introduce artificial values into flow-level measurement features, potentially contaminating the baseline dataset against which noise-injection experiments are benchmarked.
\end{enumerate}

After deletion the dataset retained all 16\,173\,208 instances that carried complete feature vectors, and no further missing-data handling was necessary.

% ============================================================
\newpage
\section{Feature Relevance Analysis PT-1 (and Removal)}
\label{sec:feature-relevance}
% ============================================================

A fundamental step before introducing controlled noise is to understand the informative content of each feature and the relationships among features. This section describes the sequence of analyses performed to characterize the feature space, identify redundancies, understand distributional properties, and quantify each feature's relationship with the target variable.

Initially this analysis was done done in order to remove from the dataset all redundant features and and features that had zero variance. Yet in chapter \ref{chap:5-Empiric-analysis-models} there will be dedicated sections to both:
\begin{itemize}
    \item \red{Experiment A:} \textbf{Experiments that used this reduced dataset}: starting from a dataset lacking any redundant feature (features with high levels of correlations) and any zero-variance feature. Those experiments were done to the purpose of finding some noise injection methods that when applied to this dataset could reliably improve performance, leading to proof of concept that it is possible to improve IDSs systems based on machine learning by injecting noise into orginal data. 
    \item \red{Experiment B:} \textbf{Experiments that used full dateset}: countrary to what was done previously, another (bigger) batch of experiments was performed starting from the full dataset. Those experiments are dedicated to embrace a more theoretical aspect of this thesis. This experiments final goal is to provide an in-depth study on how our selection of models reacts to noise input via the \texttt{Pucktrick} library.
\end{itemize}

This ramification will, obviously, lead to different conclusions in Chapter \ref{chap:6-conclusions}

% ------------------------------------------------------------
\subsection{Bivariate Correlation Analysis and Redundant Features}
\label{sec:corr-analysis}
% ------------------------------------------------------------

High pairwise correlations among features indicate redundancy: if two features carry essentially the same information, retaining both inflates computational cost without increasing the model's discriminative capacity.  Moreover, redundant features can interact adversely with regularization mechanisms and complicate the interpretation of noise-injection effects (corrupting one feature in a highly correlated pair may partially mask the corruption through the surviving redundant feature).

To systematically identify redundant features, a Pearson correlation matrix was computed over all 79 numeric features using Apache Spark's distributed \texttt{Correlation.corr()} function.  The result is a $79 \times 79$ symmetric matrix whose entry $(i,j)$ is the sample Pearson correlation coefficient between features $i$ and $j$.  Figure \ref{fig:correlation-heatmap} renders this matrix
as a heatmap using a diverging colour scale centred at zero.

\begin{figure}[H]
  \centering
  \includegraphics[width=\textwidth]{images/feature-correlation-heatmap.png}
  \caption{Pearson correlation heatmap for all 79 numeric features of the
    CIC-IDS2018 dataset.  Warm colours indicate strong positive correlation,
    cool colours indicate strong negative correlation, grey indicates
    near-zero correlation, white indicates columns with zero variance.}
  \label{fig:correlation-heatmap}
\end{figure}

The heatmap reveals pronounced block-diagonal structures, indicative of groups
of semantically related features (in the top-left-corner of the image) that are strongly inter-correlated within their group.

A threshold of $|r| \geq 0.85$ was adopted as the criterion for declaring a feature pair redundant, following the practice common in feature selection for IDS datasets.  Under this criterion,
\textbf{46 feature pairs} were identified as highly correlated.  Table~\ref{tab:high-corr-pairs-examples} lists a representative sample.

\begin{table}[htbp]
  \centering
  \caption{Representative sample of feature pairs with $|r| \geq 0.85$.}
  \label{tab:high-corr-pairs-examples}
  \small
  \begin{tabular}{llr}
    \textbf{Feature A} & \textbf{Feature B} & \textbf{$r$} \\
    Flow Duration      & Fwd IAT Tot        & 1.000 \\
    TotLen Bwd Pkts    & Subflow Bwd Byts   & 1.000 \\
    Fwd Header Len     & Bwd Header Len     & 0.999 \\
    Pkt Len Mean       & Fwd Pkt Len Mean   & 0.999 \\
    Tot Fwd Pkts       & Subflow Fwd Pkts   & 1.000 \\
    TotLen Fwd Pkts    & Subflow Fwd Byts   & 1.000 \\
    Flow IAT Std       & Fwd IAT Std        & 0.991 \\
    Bwd Seg Size Avg   & Bwd Pkt Len Mean   & 1.000 \\
    Fwd IAT Min        & Flow IAT Min       & 0.961 \\
    Fwd Act Data Pkts  & Tot Fwd Pkts       & 0.967 \\
  \end{tabular}
\end{table}

From each correlated pair, the feature deemed to carry less distinctive semantic meaning (e.g.\ sub-flow aggregates that duplicate full-flow counts, or header-length terms that duplicate packet-length terms) was designated as a removal candidate (for the \textbf{Experiments that used reduced dataset}). This process yielded \textbf{27 candidate features for removal} for \red{Experiment A} due to high inter-feature correlation.

% ------------------------------------------------------------
\newpage
\subsection{Zero-Variance Features}
\label{sec:zero-variance}
% ------------------------------------------------------------

A separate screening pass computed the standard deviation of every feature over the entire dataset.  Features with a standard deviation of exactly zero are \emph{constant} (they take the same value for every flow) and therefore provide no discriminative information whatsoever.  

Those were completely removed for \red{Experiment A} as they do not add any informative value to the dataset: on the contrary, they increase its volume, increasing computation times without providing any benefit. However (as previously specified). On the countrary, these features were retained for \red{Experiment B} to study how inserting noise into this type of feature could impact models performances.

The screening identified a total of \textbf{8 zero-variance features}:

\begin{itemize}
  \item \texttt{Bwd Blk Rate Avg}
  \item \texttt{Bwd Byts/b Avg}
  \item \texttt{Bwd PSH Flags}
  \item \texttt{Bwd Pkts/b Avg}
  \item \texttt{Bwd URG Flags}
  \item \texttt{Fwd Blk Rate Avg}
  \item \texttt{Fwd Byts/b Avg}
  \item \texttt{Fwd Pkts/b Avg}
\end{itemize}

All eight features are CICFlowMeter bulk-transfer statistics that are uniformly
zero in this capture environment, presumably because the network configuration
or the flow-generator version used did not populate these fields.  Their removal
is unconditional.

% ------------------------------------------------------------
\subsection{Reduced Feature Set}
\label{sec:reduced-feature-set}
% ------------------------------------------------------------

After removing the 8 zero-variance features and the 27 features flagged as redundant by the high-correlation criterion, the starting dataset for \red{Experiment-A} retains \textbf{43 numeric features} (\textbf{plus one date column}) and then two label columns. While, as mentione d before, for \red{Experiment-B} all \textbf{79 features} were retained.


% ============================================================
\newpage
\section{Feature Relevance Analysis PT-2}
\label{sec:feature-relevance-PT2}
% ============================================================

Following analysis were performed for both experiments. Yet for \red{Experiment A} those were applied only on the 44 remaining features, while for \red{Experiment B} those were applied on the full 79-features dataset. In following chapters all of the images will refer to the 79-features dataset as it also includes the information about the restricted dataset, but the descriptions may vary in two different sections (one for each experiment) if necessary.


% ------------------------------------------------------------
\subsection{Feature Type Classification}
\label{sec:feature-types}
% ------------------------------------------------------------

Not all features have the same measurement scale, and this
distinction is directly relevant to PuckTrick's noise-injection modules.
Gaussian noise, appropriate for continuous-valued features, is semantically
inappropriate when applied to a binary or low-cardinality categorical feature,
just as label-shuffling corruption is ill-suited to a continuous measurement.
To ensure that noise strategies are applied correctly, each feature was
classified by its \emph{effective cardinality} (the number of distinct observed
values):

\begin{itemize}
  \item \textbf{Binary} (cardinality $= 2$): flag counters that are either on
        or off for a given flow, e.g.\  \texttt{SYN Flag Cnt},
        \texttt{FIN Flag Cnt}.
  \item \textbf{Categorical} (cardinality between 3 and 200): low-cardinality
        features such as \texttt{Protocol} (three distinct values: TCP, UDP,
        ICMP encoded as integers) and window-size step features.
  \item \textbf{Continuous} (cardinality $> 200$): all high-resolution
        measurement features (durations, byte counts, inter-arrival times,
        rates, etc.).
\end{itemize}

Figure~\ref{fig:categorical-features} presents the cardinality distribution
across features.

\begin{figure}[H]
  \centering
  \includegraphics[width=\textwidth]{images/categorical_binary_features_analysis.png}
  \caption{Feature type classification for the full 79-feature set.  The bar chart (left) shows the number of distinct values per feature for low-cardinality features; the pie chart (right) summaries the proportion of binary, categorical, and continuous features.}
  \label{fig:categorical-features}
\end{figure}

The final breakdown consist of:
\begin{itemize}
    \item For the full dataset:
        \begin{itemize}
          \item \textbf{11 binary} features;
          \item \textbf{3 categorical} (low-cardinality) features;
          \item \textbf{64 continuous} features;
          \item \textbf{1 date} feature;
        \end{itemize}
    \item For the restricted dataset:
        \begin{itemize}
          \item \textbf{7 binary} features;
          \item \textbf{2 categorical} (low-cardinality) features;
          \item \textbf{34 continuous} features;
          \item \textbf{1 date} feature;
        \end{itemize}
\end{itemize}



% ------------------------------------------------------------
\subsection{Feature Fragility Analysis: Skewness and Kurtosis}
\label{sec:skew-kurtosis}
% ------------------------------------------------------------

Understanding the distributional shape of each continuous feature is essential for predicting how noise injection will affect data quality. Two distributional moments are particularly informative in this context:

\begin{itemize}
  \item \textbf{Skewness} measures the asymmetry of the distribution.  A skewness close to zero indicates an approximately symmetric distribution; large positive values indicate a long right tail (most values are small but a few are extremely large); large negative values indicate the opposite.  For noise injection, high positive skewness is especially problematic: adding zero-mean Gaussian noise to a right-skewed feature can displace values below zero (a physically meaningless outcome for quantities such as packet counts or byte totals that are strictly non-negative).
  \item \textbf{Kurtosis} (excess kurtosis, relative to the Gaussian baseline of 0) measures the heaviness of the distribution's tails.  High positive kurtosis implies that extreme values (outliers) are far more frequent than in a normal distribution.  Injecting Gaussian noise into a heavy-tailed feature risks corrupting exactly those extreme values that are most discriminative for detecting volumetric attacks (which manifest precisely as statistical outliers in byte-count and rate features).
\end{itemize}

Those descriptive statistics were computed for all numeric features using Spark's distributed aggregation functions.  Figure~\ref{fig:skew-kurtosis} displays skewness and kurtosis as side-by-side horizontal bar charts at following page:

\begin{landscape}
\begin{figure}[H]
    \centering
    \hspace*{-1.2cm}\makebox[\textwidth][c]{\includegraphics[width=1.35\linewidth]{images/feature-skew-kurtosis.png}}
    \caption{Skewness (left) and excess kurtosis (right) for the numeric
    features. Red bars indicate values exceeding the fragility threshold
    ($|\text{skewness}| > 2$ and kurtosis $> 7$, respectively), signaling
    features likely to be significantly affected by additive Gaussian noise.}
    \label{fig:skew-kurtosis}
\end{figure}
\end{landscape}

The analysis reveals an extreme distributional asymmetry prevalent across the
feature set.  Notable examples include:

\begin{itemize}
  \item \texttt{Pkt Len Var}: skewness $\approx 1863.6$, kurtosis $\approx 3\,600\,000$ , in the most extreme case;
  \item \texttt{TotLen Fwd Pkts}: skewness $\approx 836.7$, indicating a distribution dominated by very small values (short flows) with an extended right tail of large data transfers;
  \item \texttt{Tot Bwd Pkts}: skewness $\approx 170.8$;
  \item \texttt{Flow Byts/s} and \texttt{Flow Pkts/s}: these rate features exhibit numerically infinite values due to zero-duration flows (a division-by-zero artifact of the flow meter), which must be handled carefully prior to any noise injection.
\end{itemize}

% ............................................................
\subsubsection{Feature Fragility Classification}
\label{sec:fragility}
% ............................................................

Based on the joint evaluation of skewness and kurtosis, each continuous feature was assigned a \emph{fragility level} indicating its susceptibility to statistical distortion by PuckTrick's noise-injection procedures:

\begin{itemize}
  \item \textbf{High fragility} (high skewness \emph{and} high kurtosis): features whose distributions are both asymmetric and heavy-tailed. Gaussian noise will simultaneously generate invalid negative values and disrupt the discriminative outlier structure.
  \item \textbf{Medium fragility} (high skewness \emph{or} high kurtosis but not both): features with one adverse distributional property but not the other.
  \item \textbf{Low fragility} (robust): features whose distributions are approximately symmetric and have near-Gaussian tails; Gaussian noise injection is both statistically appropriate and unlikely to produce artefactual values.
\end{itemize}

In total, \textbf{57 out of 79 features} results as higly subjective ot noise, classifying it as high fragility. While only 6 presents medium fragility and 14 presents low fragility.

Representative examples from each fragility class are shown in Table \ref{tab:fragility-examples}.

\begin{table}[H]
  \centering
  \caption{Representative features from each fragility class.}
  \label{tab:fragility-examples}
  \small
  \begin{tabular}{llrrl}
    \textbf{Feature} & \textbf{Fragility} & \textbf{Skewness} & \textbf{Kurtosis} & \textbf{Notes} \\
    Pkt Len Var            & High    & 1863.6  & 4\,295\,779 & Extreme dominance of small values \\
    TotLen Fwd Pkts        & High    & 836.7   & 1\,808\,721 & Right-tail dominated \\
    Tot Fwd Pkts           & High    & 88.6    & 8\,271      & High-cardinality count \\
    Bwd Pkts/s             & High    & 9.07    & 103.71      & Rate feature w/ inf. values \\
    Flow Duration          & Medium  & 2.64    & 5.53        & Moderate tail weight \\
    Bwd Pkt Len Max        & Medium  & 1.67    & 55.25       & Moderate tail weight \\
    Fwd Seg Size Min       & Low     & 0.26    & 0.56        & Near-Gaussian \\
  \end{tabular}
\end{table}

This fragility map constitutes a theoretically grounded checklist for interpreting the downstream experimental results: when a noise-injection run targets high-fragility features, a larger degradation in model performance is expected than when robust features are perturbed.

% ------------------------------------------------------------
\newpage
\subsection{Feature–Target Correlation Analysis}
\label{sec:feat-target-corr}
% ------------------------------------------------------------

Having characterized the intra-feature structure and distributional properties, the final step of the dataset analysis quantifies the relationship between each feature and the target variable. This is done to add a third index to fragility measures: having features that are already evaluated being fragile while also being highly correlated to a target column may allow to find some of the most fragiles features of the entire Dataset. 

Two complementary statistical measures were employed, one for each level of the label hierarchy:

\begin{enumerate}
  \item \textbf{Point-biserial correlation with \texttt{Label\_generic} (for binary target):} the \textit{Pearson correlation} coefficient between a continuous feature and the binary (\texttt{Malign} / \texttt{Benign}) encoded target. This measures how much the feature's mean differs between benign and malign flows, normalized by the overall variability.
  
  \item \textbf{ANOVA F-statistic with \texttt{Label} (for multi-class target):} the one-way analysis of variance F-statistic measures the ratio of between-class variance to within-class variance. A high F-statistic indicates that a feature takes substantially different mean values across the individual attack categories, making it useful for fine-grained attack classification.
\end{enumerate}

Figure \ref{fig:target-correlation} presents the two rankings side by side, colour-coded by the fragility class defined in
Section~\ref{sec:fragility}.

\begin{figure}[htbp]
  \centering
  \includegraphics[width=\textwidth]{images/fragile-features-target-correlation.png}
  \caption{Feature relevance with respect to the target variable.
    \emph{Left:} features ranked by $|\rho|$ (absolute Pearson correlation)
    with the binary \texttt{Label\_generic}.
    \emph{Right:} features ranked by the ANOVA F-statistic with the
    multi-class \texttt{Label} (log scale).
    Colour coding: red = high fragility, orange = medium fragility,
    green = low fragility.}
  \label{fig:target-correlation}
\end{figure}

\paragraph{Correlation with the binary label.}
The feature exhibiting the highest absolute correlation with \texttt{Label\_generic} is \texttt{Fwd\_Seg\_Size\_Min} ($|\rho| \approx 0.43$). Crucially, the \emph{maximum} observed correlation is approximately below $50\%$. Overall this is a modest value that places an inherent limit on how well any linear or correlation-based model can separate the two classes using a single feature. 
This is a direct consequence of the absence of payload-level information: since
the dataset contains only transport-layer statistics measured on encrypted or
high-volume traffic, no individual feature provides a sharp decision boundary for the most generic task. This means that for our experiments (Both for \red{Experiment A} and \red{Experiment B}) there will most likely not be any huge downgrades when performing noise injection on this binary label. 

\paragraph{Discrimination across attack types (ANOVA).}
The ANOVA analysis paints a complementary picture.  Features such as \texttt{Tot\_Fwd\_Pkts} (F $\approx 2.0 \times 10^6$), \texttt{Bwd Pkts/s} (F $\approx 1.8 \times 10^6$), and \texttt{TotLen Fwd Pkts} (F $\approx 7.4 \times 10^5$) exhibit extremely high between-class variance relative to within-class variance, indicating that they are excellent discriminators \emph{across attack categories}, even though their Pearson correlations with the binary label are moderate.  This asymmetry arises because different attack types produce qualitatively different flow profiles (e.g., DoS floods generate enormous packet counts while brute-force attacks generate numerous short flows), yet all attack flows collectively differ from benign flows in a more diffuse, multi-dimensional fashion.

\paragraph{Fragility vs.\ relevance tension.}
A particularly significant observation emerging from Figure \ref{fig:target-correlation} is that the most discriminative features tend to fall in the \emph{high-fragility} category (red bars).  This creates a direct tension with noise injection: the features that contribute most to classifier performance are precisely the features whose distributions are most vulnerable to distortion by Gaussian noise.  This tension is one of the central experimental questions addressed in Chapter~\ref{chap:Empiric-analysis-models}: "can deep learning models maintain their discriminative capability when their most informative-yet-fragile features are corrupted by PuckTrick?"

% ============================================================
\section{Summary}
\label{sec:dataset-summary}
% ============================================================

Table~\ref{tab:dataset-summary} provides a consolidated overview of the
CIC-IDS2018 dataset properties as characterised in this chapter.

\begin{table}[htbp]
  \centering
  \caption{Summary of CIC-IDS2018 dataset characteristics.}
  \label{tab:dataset-summary}
  \begin{tabular}{ll}
    \toprule
    \textbf{Property} & \textbf{Value} \\
    \midrule
    Total instances               & 16\,173\,208 \\
    Original features             & 79 \\
    Label columns                 & 2 (\texttt{Label}, \texttt{Label\_generic}) \\
    Zero-variance features removed & 8 \\
    High-correlation features removed & 27 \\
    Retained features for \red{Experiment A}  & \textbf{44} \\
    \midrule
    Feature types (for \red{Experiment A}) & 34 continuous, 7 binary, 2 categorical, 1 date \\
    Feature types (for \red{Experiment B}) & 64 continuous, 11 binary, 3 categorical, 1 date \\
    Class distribution            & Severely imbalanced (Benign $\gg$ Malign) \\
    \midrule
    High-fragility features       & Vast Majority for both datasets ($\approx 70\%$)\\
    Medium-fragility features     & Most rare features (($\approx 10\%$)) \\
    Low-fragility features        & Somewhat present features ($\approx 20\%$) \\
    \midrule
    Max $|\rho|$ with binary label & $\approx 0.43$ (\texttt{Fwd Seg Size Min}) \\
    Max ANOVA F with multi-class label & $\approx 2.0\times10^6$ (\texttt{Tot Fwd Pkts}) \\
    \bottomrule
  \end{tabular}
\end{table}


The analysis presented in this chapter establishes the informational context
necessary to interpret the noise-injection experiments.  The following points
are especially relevant for subsequent chapters:

\begin{enumerate}
  \item The \textbf{severe class imbalance} requires that model performance be
        evaluated with metrics beyond accuracy (e.g., F1-score,
        precision-recall curves) and that the impact of noise be assessed
        class-conditionally.
  \item The \textbf{extreme distributional skewness} of most features implies
        that PuckTrick's Gaussian noise module will, in many cases, violate
        the non-negativity constraints of the measurement domain.  This makes
        the observed performance degradation attributable not only to the
        statistical displacement of feature values but also to the creation of
        physically invalid instances.
  \item The \textbf{fragility–relevance tension} (the fact that the most
        predictive features are also the most fragile) means that even
        moderate noise magnitudes can have disproportionate effects on model
        performance, a hypothesis central to the experimental design.
  \item The \textbf{temporal signal} embedded in the \texttt{Timestamp}
        column provides contextual information that complements the
        flow-statistics features.  Its corruption under PuckTrick's label
        inconsistency module constitutes a distinct and interesting experimental
        scenario.
\end{enumerate}

Taken together, these properties render the CIC-IDS2018 dataset a demanding and
realistic testbed for evaluating the effects of controlled data corruption on
deep learning-based intrusion detection systems.


\end{document}